\begin{abstract}
Multimodal visual object tracking augments an RGB camera with auxiliary sensors (e.g., thermal, depth, or event) to maintain robustness in adverse conditions. However, real-world sensor rigs frequently suffer from de-synchronization or signal dropout, rendering the auxiliary modality intermittently missing. Conventional multimodal trackers, trained on complete modality pairs, degrade sharply under these conditions. While recent robust trackers mitigate this by adaptively scaling their capacity or masking the absent stream, they fundamentally treat the problem as \emph{passive coping} rather than \emph{active compensation}, leaving the cross-modal correlation unused. In this paper, we propose FlexTrack-V2, a unified tracking framework that actively hallucinates missing information. Our core insight is that the present modality carries a strong prior about the absent one. To exploit this, we introduce the Bilateral Modality-specific feature Reconstruction (BMR) module, which couples with a Heterogeneous Mixture-of-Experts (HMoE) to actively synthesize missing features and dynamically allocate computation. To train this generative pathway stably without an external teacher, we propose Curriculum Missing Augmentation (CMA) with online self-distillation, which anneals the missing rate and aligns degraded predictions with complete-modality representations. Built on a shared Fast-iTPN encoder and a Mamba interaction neck, FlexTrack-V2 uses a single set of parameters for RGB-Thermal, RGB-Depth, and RGB-Event tracking. Extensive experiments across nine benchmarks demonstrate that FlexTrack-V2 achieves parity with state-of-the-art models under complete modalities and yields substantial, consistent improvements under missing-modality regimes.
\end{abstract}

\begin{IEEEkeywords}
Multimodal Object Tracking, Missing Modalities, Mixture-of-Experts, Feature Reconstruction.
\end{IEEEkeywords}
